\documentclass[a4paper, 11pt]{report}
\usepackage[utf8]{inputenc}
\usepackage[T1]{fontenc}
\usepackage[french]{babel}

\begin{document}
\noindent
{\bf UNIVERSITÉ DE LOMÉ \hfill ANNÉE UNIVERSITAIRE \\
FACULTÉ DES SCIENCES \hfill 2020-2021 }

\begin{center}
  {\bf TD de INF203 Algorithmes et Programmation Fortran}
\end{center}

\subsubsection{\underline{Exercice 1}}

\noindent
L'énergie potentielle inermoléculaire peut être calculée par la formule de Renards-jones: 
$$ E(R) = 4\varepsilon((\sigma/R)^2 - (\sigma/R)^6) $$
où $\varepsilon$ et $\sigma$ sont les paramètres caractéristiques des molécules mises en jeu. \\

\vspace{0,5 cm}

\noindent
Écrire un programme permettant de calculer $E(R)$ pour des valeurs de $R$ comprises entre 1 et 5 \AA \hspace{0,05 cm} par pas de 0,25, les valeurs des paramètres $\varepsilon$ et $\sigma$ sont donnés dans le tableau ci-contre 



\begin{table}[h]
  \begin{flushright}
  \begin{tabular}{||l c r||}
    \hline
    \hline
    Atomes & $\varepsilon(kcla/mole)$ & $\sigma($\AA$)$ \\
    \hline
    \hline
    He & 0,12 & 2,6 \\
    Ne & 0,07 & 2,8 \\
    Ar & 0,24 & 3,4 \\ 
    Kr & 0,35 & 3,8 \\
    Xe & 0,45 & 4,0 \\
    \hline
    \hline
  \end{tabular}
  \end{flushright}
\end{table}

\noindent
Enregistre dans des fichiers (un fichier par élément gaz rare) les données dans le but de reproduirele graphe de l'énergie potentielle en foction de la position $R$. \\
Chaque fichier doit contenir deux colonnes correspondant aux valeurs succesives de $R$ et $E(R)$

\subsubsection{\underline{Exercice 2}}

\noindent
L'énergie potentielle d'une particule est donnée par la relation $ E = \frac{1}{r^3} - \frac{1}{r^2} $ .

\begin{enumerate}
  \item[(a)] Écrire un programme avec une boucle \verb|DO| permettant de calculer $E$ pour des valeurs de $r$ entières comprise entre 1 et 10. Le résultat sera enregistré dans un fichier.
  \item[(b)] Idem avec $r$ variant de 0,5 et 5 avec un pas de 0,5.
  \item[(c)] Idem avec $r$ variant de 0,7 à $X$ avec un pas de 0,3. La valeur de $X$ devrait être donnée à l'exécution.
\end{enumerate}

\subsubsection{\underline{Problème}}

\begin{enumerate}
  \item Écrire une \verb|function| externe dénomée \verb|fact(N)| permettant de calculer le factoriel d'un nombre entier positif $N$ . On rappelle que 0! = 1 .
  \item On veut calculer $e^x$ à l'aide d'un développemnt limité : 
    $$ e^x = 1 + \frac{x}{1!} + \frac{x^2}{2!} + \cdots + \frac{x^i}{i!} = a_0 + a_1 + a_2 + \cdots a_i $$ . \\
    On décide d'arrêter le développment lorsque le rapport du terme $a_i$ d'ordre $i$ sur la somme de tous les précédents est inférieur à une valeur donnée $\varepsilon $ . \\ 
    
    \vspace{0,125 cm}
    \noindent
    En se servant de la \verb|function| externe \verb|fact(N)| calculant les factoriels, écrire une \verb|function| externe \verb|exp(x)| approchant l'exponentiel d'une valeur réelle $x$ . \\
    On remarquera que chaque terme de cette somme peut être obtenu à partir du précédent par :
    $$ a_k = a_{k - 1} * \frac{x}{k} $$ ..
\end{enumerate}
\end{document}